% RESEARCH_DRAFT_SHA256: e04ba623816656907e9c914cf767ddec0a35b8b36d202c39097413ac2c6a90fd
% Layout and source-reference wrappers expanded for the web renderer.
\documentclass[11pt]{article}
\usepackage[margin=25mm]{geometry}
\usepackage{amsthm}
\usepackage{fontspec,unicode-math}
\setmainfont{Latin Modern Roman}
\setmathfont{Latin Modern Math}
\usepackage{amsmath,amssymb,mathtools,xurl,hyperref}
\hypersetup{colorlinks=true,urlcolor=blue,linkcolor=blue}
\setcounter{secnumdepth}{0}
\setlength{\parindent}{0pt}
\setlength{\parskip}{5pt}
\setlength{\emergencystretch}{3em}
\newcommand{\sref}[2]{\hyperlink{source:#1:#2}{[S#2]}}
\newcommand{\cataloguescope}[1]{\paragraph{Source record.}#1}
\newcommand{\floor}[1]{\left\lfloor#1\right\rfloor}
\newtheorem{proposition}{Proposition}
\newtheorem{lemma}[proposition]{Lemma}
\begin{document}
\paragraph{Research manuscript.} OpenAI \texttt{gpt-6-astra}, reasoning \texttt{ultra}; 11 October 2026.
\paragraph{Preserved source body.} The complete catalogue mathematical text and sources follow with layout and source links expanded, including source qualifications and its historical preparation notes. New research begins at the section headed Research attempt.

% UnsolvedMath ID 20002201; source catalogue code AIM-INFRASTRUCTURE-0101
\section{A mobility-and-care adjusted compensation and incidence framework}\label{20002201}
{\small UnsolvedMath ID 20002201 · Computer Science}
\subsection{Definitions and mathematical statement}

\paragraph{Shared notation.}Standard number systems are $\mathbb N,\mathbb Z,\mathbb Q,\mathbb R,\mathbb C$. Specialized notation and index ranges follow the cited source and the definitions below.


Structural employment conditions concern recurring features of positions and institutions. A pay comparison requires defined compensation measures and comparable populations. Relocation, caregiving expenses, and healthcare access involve different resources and outcomes; combining them does not define a single mathematical quantity.
\paragraph{Record scope and qualifications.}
The catalogue title is retained as metadata. Named certificates, contracts, bounds or frameworks in that title are not asserted as results of the original source extract.
The source gives no jurisdiction, institution, intervention, or quantified evaluation criterion. It is a social and institutional agenda; no legal or medical advice is supplied.

Catalogue screening limitation: Community, history, equity, or organizational policy discussion, rather than a mathematical statement.

The source snapshot labels this record partially solved; this is catalogue metadata rather than independent verification of a mathematical result or the current completion of a project.
\paragraph{Statement recorded in the supplied source.}
Less Actionable But Nonetheless Important Issues

How do we deal with the structural problem of expecting mathematicians to move every 2 years for post-docs? How can we work to change the gender pay gap? How can we financially support mathematicians who would like to have kids, especially those who have kids in ways that are financially draining? How can we provide better access to healthcare for trans mathematicians? \href{http://aimpl.org/gemscombin/7/}{[S2]}
\paragraph{Source status.}
The supplied catalogue labels this source record partially solved; that metadata was checked against the supplied snapshot on 2026-10-11. The label is preserved for provenance and is not an independent verification of present open status or mathematical correctness.
\subsection{Short English statement}
Discuss employment mobility, pay, caregiving, and healthcare barriers affecting mathematicians.
\subsection{Sources}
\begin{itemize}
\item[S1] Ulam AI and the UnsolvedMath Contributors, UnsolvedMath: A Curated Collection of Open Mathematics Problems, record 20002201 (AIM-INFRASTRUCTURE-0101). Catalogue curation and metadata: CC BY 4.0. The source extract is retained; definitions, English summaries and qualifications are editorial additions.
\url{https://huggingface.co/datasets/ulamai/UnsolvedMath}.
\item[S2] AIM Workshop Problem Lists, source linked by catalogue record AIM-INFRASTRUCTURE-0101.
\url{http://aimpl.org/gemscombin/7/}.
\end{itemize}
\paragraph{Preparation and review.}
The original source and UnsolvedMath contributors retain attribution for the source extract. Codex (GPT-6; exact serving revision not exposed) assisted with definitions, English summaries, scope notes and formatting. No new human mathematical authorship or independently verified solution is asserted. Permanent source ID 20002201 and code AIM-INFRASTRUCTURE-0101 checked against the complete supplied registry and current repository. Distinct upstream records, including repeated questions and broad or damaged extracts, are retained for complete coverage of the source CS category. A repeated record is not a new independent conjecture. Complete mathematical-equivalence and current-literature review has not been performed. Preparation and final review follow contribution prompts C and D, with source gaps disclosed. Structural and rendering checks do not establish a new theorem or open status.
\label{end:20002201}

\section{Research attempt}
\subsection{Literature, source context and scope}
The complete catalogue body above retains the original agenda, source status, and historical preparation. The AIMPL problem page could not be retrieved. The supplied complete source transcription was read, and the accessible primary AIM workshop announcement was checked for context. The announcement is not a substitute for verification of the inaccessible problem page. The catalogue's earlier proposed frameworks and research summaries are context, not established solutions credited to this attempt. The following is an explicitly limited investigation of one route; it neither substitutes a formalized toy model for the entire source agenda nor claims empirical or deployed-software results. Elementary auxiliary facts are not new original-target progress.

\subsection{A restricted resource-allocation route}
The source combines mobility, compensation, family support and healthcare access. A mathematical allocation model cannot resolve those institutional questions. To test one limited way of allocating a fixed support budget, let $c_1,\ldots,c_n\ge0$ be externally supplied commensurable burdens and let $0\le B\le\sum_i c_i$. A transfer $f_i\ge0$ uses budget $\sum_i f_i\le B$ and leaves residual burden $(c_i-f_i)_+$, where $u_+=\max(u,0)$.
\begin{proposition}
For $0<B<\sum_i c_i$, there is a unique $t\in(0,\max_i c_i)$ with $\sum_i(c_i-t)_+=B$. Transfers $f_i=(c_i-t)_+$ minimize the largest residual burden, with optimal value $t$. For $B=0$ the optimum is $\max_i c_i$, and for $B=\sum_i c_i$ it is zero.
\end{proposition}
\begin{proof}
The function $F(t)=\sum_i(c_i-t)_+$ is continuous, with $F(0)=\sum_i c_i$ and $F(\max c_i)=0$. It is strictly decreasing on $[0,\max c_i]$ if the maximum is positive: for any two distinct points below the maximum, at least one term strictly decreases, while none increases. The intermediate value theorem gives the unique interior root for the stated interior budgets. The proposed transfers sum to $B$ and leave each residual equal to $\min(c_i,t)$, so their maximum is $t$. If some allocation had maximum residual at most $r<t$, it would need $f_i\ge(c_i-r)_+$ for every $i$, hence total at least $F(r)>F(t)=B$, a contradiction. The endpoint budgets give the asserted optima by nonnegativity and the allocations $f=0$ or $f=c$. If all burdens are zero, both endpoint conclusions coincide.
\end{proof}
For $c=(8,4,2)$ and $B=6$, the equation gives $t=3$ and $f=(5,1,0)$; residuals are $(3,3,2)$. Equal transfers $(2,2,2)$ instead leave maximum 6. This illustrates the model's objective, not an assertion that minimizing the maximum of these invented numbers is the community's chosen fairness criterion.
\subsection{Verification and remaining gaps}
The example and all integer allocations with total at most 6 were enumerated; the smallest maximum residual was 3. The proof covers real transfers as well. No burdens were measured, budget proposed, financial or healthcare recommendation made, or structural employment reform evaluated. Noncommensurable needs, unequal transfer effectiveness and eligibility choices are outside the model and material to the original agenda.

\section{References and source attribution}
Primary source: \url{http://aimpl.org/gemscombin/7/}. Attempted primary AIMPL page retrieval; failed with timeout or gateway error. The supplied full source transcription was read; the inaccessible live page was not independently inspected.

\url{https://aimath.org/pastworkshops/gemscombin.html}. Read the primary March 2023 workshop announcement and its community inclusion agenda.

All retrievals were checked on 11 October 2026. Original contributor and catalogue attribution remains in the preserved Sources section. The elementary arguments below are self-contained; their exposition is not a claim of first discovery.


\paragraph{Integration review.} Codex (GPT-6; exact serving revision not exposed) reviewed the written arguments and expanded layout and source-reference wrappers. Repository metadata and the production web renderer were checked using Python 3.13 and JavaScript with MathJax 3.2.2. This model-assisted review is not independent human proof verification; no PDF compilation or formal proof checking is asserted.

\section{Completion Estimate}
\noindent\textbf{COMPLETION: 0\%}
No substantiated new progress toward the entire original target is claimed. The worked results and failed route are delimited above; a known result, illustrative model, or auxiliary obstruction is not counted as original-target completion.

\section{Final status and attribution}
\textbf{Attempt status: unresolved.} This describes this research attempt, not a new ruling on the source catalogue's status or on external literature claims.
The mathematics and this self-audit were generated by OpenAI \texttt{gpt-6-astra}, reasoning \texttt{ultra}, on 11 October 2026. Codex (GPT-6; exact serving revision not exposed) orchestrated the batch and contributed formatting/repository review. No human mathematical authorship, independent human proof checking, or formal verification is claimed. The historical preparation paragraphs in the reproduced catalogue body concern that body only.

Contribution prompt A was applied to the complete definition and supplied source record, followed by the adversarial checks of prompt B. Local reproduction files are \texttt{batch013/checks.py}, \texttt{computations.json}, and the per-ID audit JSON. They were not public links at preparation time. Numerical checks are finite and do not prove unrestricted results. No LaTeX engine was found locally; compilation was not performed by this research agent.

\end{document}
